\section{Attack Suite and Configuration}
% TODO[Kiệt]:
% Goal:
% - Make this section the source of truth for attack names and configurations.
% Flow:
% - Separate source/training attacks from evaluation-only attacks and AutoAttack.
% - Define PGD-$\ell_\infty$ once as the 20-step cross-entropy source attack.
% Include:
% - Ensure tables cover attack, norm, budget, steps, loss/objective, role, and whether used for training/evaluation.
% - Cite docs/05_notation_registry.md names in formal table/caption text.
% Avoid:
% - Repeating code-style names like `PGD20-CE` in prose.
% Target length: 4-6 paragraphs plus configuration tables.
Stage 2 evaluates whether robustness learned from a small set of source attacks generalizes to a broader attack suite. The attack suite is divided into two groups: source attacks used during training and evaluation attacks used only for validation or final testing.

\subsection{Training attacks}
The main training setup uses two source attacks:
\[
    \mathcal{A}_{\mathrm{src}}
    =
    \{\text{PGD-}\ell_\infty,\text{DDN-}\ell_2\}.
\]
PGD-\(\ell_\infty\) represents the standard first-order \(\ell_\infty\) adversarial
training setting. It uses cross-entropy loss and perturbation budget
\(\varepsilon_\infty = 8/255\). DDN-\(\ell_2\) is based on the Decoupled Direction and
Norm attack of Rony et al.~\cite{rony2019decoupling}. DDN is an efficient gradient-based
\(\ell_2\) attack that decouples the perturbation direction from its norm: the direction
is updated using gradient information, while the perturbation norm is adapted depending
on whether the current example is already adversarial.



Table~\ref{tab:training_attacks} summarizes the source attacks used during training.
Both attacks are used to generate mixed adversarial minibatches for Multi-ATtack
ERM, AttackDRO, AttackDRO++, and the proposed AttackDRO++ Anchor35 GradFP method.

\begin{table}[H]
\centering
\caption{Training attack configuration.}
\label{tab:training_attacks}
\begin{tabular}{llllll}
\toprule
Attack & Norm & Budget & Steps & Loss / objective & Role \\
\midrule
PGD-{$\ell_\infty$} & $\ell_\infty$ & $8/255$ & 20 & Cross-entropy & Source attack \\
DDN-$\ell_2$ & $\ell_2$ & $0.5$ & 20 & DDN objective & Source attack \\
\bottomrule
\end{tabular}
\end{table}

ForMulti-ATtack training, each minibatch is split across the source attacks and adversarial examples are generated using the assigned attack. The same mixed-batch source attack construction is also used by AttackDRO, AttackDRO++, and the proposed AttackDRO++ Anchor35 GradFP method. Thus, the compared methods share the same adversarial data generation process and differ mainly in how they weight or group the generated examples.

\subsection{Evaluation Attacks}
\label{sec:evaluation_attacks}

The evaluation suite includes both source attacks and held-out attacks. The two source
attacks used during training are also included in validation and test evaluation, while
the remaining attacks are held out from training. This protocol directly tests the
attacks-as-domains hypothesis: a robust model should not only perform well on the attacks
used for optimization, but should also generalize to attacks with different perturbation
geometries and optimization dynamics.


For validation and final testing, we evaluate each checkpoint using the attack suite in
Table~\ref{tab:evaluation_attacks}. The main aggregate metric, denoted Mean(8), is
computed over the eight attacks in the table. AutoAttack is reported separately and is
therefore not included in Mean(8).
\begin{table}[H]
\centering
\caption{Validation and test attacks used for Mean(8).}
\label{tab:evaluation_attacks}
\renewcommand{\arraystretch}{1.08}
\setlength{\tabcolsep}{7pt}
\begin{tabular}{lcccc}
\toprule
Attack & Norm & Budget & Steps & Group \\
\midrule
FGSM-RS              & $\ell_\infty$ & $8/255$ & 1   & Held-out $\ell_\infty$ \\
PGD-$\ell_\infty$   & $\ell_\infty$ & $8/255$ & 20  & Source \\
TPGD                 & $\ell_\infty$ & $8/255$ & 10  & Held-out $\ell_\infty$ \\
MI-FGSM              & $\ell_\infty$ & $8/255$ & 10  & Held-out $\ell_\infty$ \\
PGD-$\ell_2$         & $\ell_2$      & $0.5$   & 10  & Held-out $\ell_2$ \\
DDN-$\ell_2$         & $\ell_2$      & $0.5$   & 20  & Source \\
DeepFool-$\ell_2$    & $\ell_2$      & $0.5$   & 50  & Held-out $\ell_2$ \\
CW-$\ell_2$          & $\ell_2$      & $0.5$   & 200 & Held-out $\ell_2$ \\
\bottomrule
\end{tabular}
\end{table}

For aggregate analysis, we report four summary metrics derived from the same eight
attacks: Seen, Held-out \(\ell_\infty\), Held-out \(\ell_2\), and Mean(8). Their exact
definitions are given in Section~\ref{sec:evaluation_metrics}.

\paragraph{AutoAttack configuration}

AutoAttack is used as an additional standardized \(\ell_\infty\) stress test following
Croce and Hein~\cite{croce2020reliable}. It is not included in Mean(8), because it is
computationally more expensive and follows a separate evaluation protocol.

The AutoAttack configuration uses the standard \(\ell_\infty\) setting with perturbation
budget \(\varepsilon_\infty = 8/255\). In early experiments and ablations, AutoAttack is
sometimes evaluated on a subset of 512 test samples as a sanity check. These 512-sample
results are useful for detecting severe robustness collapse, but they are not used as the
primary ranking criterion because the sample size is smaller and the variance is higher.
For the final stress test, AutoAttack is evaluated on the full test set for selected
representative seeds.

\begin{table}[H]
\centering
\caption{AutoAttack configuration.}
\label{tab:autoattack_config}
\renewcommand{\arraystretch}{1.08}
\setlength{\tabcolsep}{10pt}
\begin{tabular}{lccc}
\toprule
Attack & Norm & Budget & Role \\
\midrule
AutoAttack & $\ell_\infty$ & $8/255$ & Separate stress test \\
\bottomrule
\end{tabular}
\end{table}
